% ID: SEL-D0070-OA-E02
% ESTADO: cotejada. Segunda pasada visual y borrador compilado revisados.
% Fuente: Selectividad/sel_2009_matematicas_aplicadas/mat apli CC SS04.pdf
% SHA-256: bdea2fa323085e0db4eacff922ba9f403f58ef4e6bb4f2b569a0b857554c9ef2
% Página: 1; opción A; ejercicio original 2.

Un almacenista de frutas ha estimado que el beneficio que le produce cada kilogramo (kg) de fresas depende del precio de venta de acuerdo con la función
\[
B(x)=-x^2+4x-3
\]
siendo $B(x)$ el beneficio por kg y $x$ el precio de cada kg, ambos expresados en euros.
\begin{enumerate}[label=\alph*)]
\item \textbf{(1.25 puntos)} \textquestiondown{}Entre qué precios se producen beneficios para el almacenista?
\item \textbf{(1.25 puntos)} \textquestiondown{}Qué precio maximiza los beneficios?
\item \textbf{(0.5 puntos)} Si tiene en el almacén 10000 kg de fresas, \textquestiondown{}cuál será el beneficio total máximo que podrá obtener?
\end{enumerate}
